\chapter{Reports, and the PDF helper that has to be sandboxed}\label{ch:pdf}

Two files produce a Test Report. \texttt{src/\allowbreak{}ate/\allowbreak{}report/\allowbreak{}html.py} renders the self-contained HTML from
the recorded run; \texttt{src/\allowbreak{}ate/\allowbreak{}report/\allowbreak{}pdf.py} turns that HTML into PDF in a separate process, and
\texttt{src/\allowbreak{}ate/\allowbreak{}report/\allowbreak{}pdf\_\allowbreak{}job.py} is the \texttt{QProcess} plumbing that runs it. A Test Report is not
optional (spec 21a), so a failure here must never crash a run.

\section{The HTML}

\texttt{write\_\allowbreak{}report(store, run\_\allowbreak{}uuid, config)} renders through the Jinja template
\texttt{src/\allowbreak{}ate/\allowbreak{}report/\allowbreak{}templates/\allowbreak{}report.html.j2} and writes under \texttt{reports.directory} as

\begin{lstlisting}[language=ATEShell]
<UUT Type>/<UUT Type>_<serial>_<local start yyyymmdd-HHMMSS>_<PASS|FAIL>.html
\end{lstlisting}

Windows-invalid characters are replaced, and an existing file is never overwritten: a free
\texttt{-2}, \texttt{-3} suffix is used instead. The name is only free when neither the HTML nor its PDF
exists, because the PDF takes the same stem. \texttt{generated\_at} freezes the footer stamp for a
byte-deterministic write, which is what makes the report tests able to compare output.

\texttt{src/\allowbreak{}ate/\allowbreak{}report/\allowbreak{}log\_\allowbreak{}export.py} is the Run log export (packet 24a): CSV, JSON Lines or text,
filtered by level, source and owner --- the same filters \texttt{ExportLogDialog} shows.

\section{Why the PDF is a separate process}

Rendering needs Qt WebEngine. Running it in the application's process would put Chromium's renderer in
the same process as the Test Run, and the render is slow: a run's \texttt{FINISHED} must not wait for it.
So the helper is \texttt{python -m ate \_\allowbreak{}pdf <html> <pdf> --paper A4|Letter [--unsandboxed]} started by
\texttt{start\_pdf\_job}, with \texttt{PDF\_\allowbreak{}TIMEOUT\_\allowbreak{}MS = 60000}. The helper is parented to the session
that started it, and \texttt{stop\_pdf\_jobs()} must run before the session dies: a \texttt{QProcess}
destroyed mid-render by Qt makes the job's \texttt{finished} callback fire against a session whose Qt
side is gone, and that aborted the test run at teardown (packet 25b review). It is why
\texttt{HeadlessRunSession.close} and \texttt{\_\_del\_\_} both call it, and why
\texttt{conftest.py} has an autouse fixture that calls it after every test.

\section{The sandbox decision, asked once}

Chromium's sandbox cannot run as root or as a Windows administrator, and turning it off silently is a
security decision nobody should make by accident. \texttt{pdf\_\allowbreak{}sandbox\_\allowbreak{}choice()} is built once at launch
and handed to the session and to History, so neither of them reads the environment or guesses at
privilege. The decision, per PDF:

\begin{itemize}
  \item not privileged --- \texttt{PDF\_SANDBOXED}, nothing is asked. \texttt{--allow-unsandboxed-pdf}
    has no effect, and the log says so;
  \item privileged and \texttt{--allow-unsandboxed-pdf} --- \texttt{PDF\_UNSANDBOXED}, no dialog, logged
    as a WARNING;
  \item privileged and there is a dialog to ask --- the Operator's answer. No means \texttt{PDF\_SKIPPED},
    and the HTML Test Report is still written;
  \item privileged with no way to ask (headless, \texttt{ate \_selftest --run}) --- \texttt{PDF\_SKIPPED}.
\end{itemize}

Two environment variables can switch WebEngine's sandbox off by inheritance,
\texttt{SANDBOX\_\allowbreak{}OFF\_\allowbreak{}VARIABLES} --- \texttt{QTWEBENGINE\_\allowbreak{}DISABLE\_\allowbreak{}SANDBOX} and
\texttt{QTWEBENGINE\_\allowbreak{}CHROMIUM\_\allowbreak{}FLAGS} with \texttt{--no-sandbox} --- and \texttt{\_\allowbreak{}strip\_\allowbreak{}inherited\_\allowbreak{}sandbox\_\allowbreak{}off}
removes them from the child's environment so a Station cannot get an unsandboxed render without choosing
it.

The test seam is \texttt{ATE\_ELEVATED} (\texttt{ELEVATION\_ENV}): \texttt{conftest.py} pins it to 0 for
every test, because a GitHub Windows runner is an administrator and a container runs as root, and a
privileged session asks a question that \texttt{exec()} would block forever on an offscreen platform. A
test that wants to be privileged sets \texttt{ATE\_ELEVATED=1}, and it reaches a spawned \texttt{ate \_pdf}
child through the inherited environment.

\section{Proving a report rendered}

A 13-byte stub PDF is enough for signal plumbing and is what \texttt{conftest.py} installs for every
run that is not specifically about the PDF. Four tests use the real helper, because ``the PDF has text in
it'' can only be proven against a PDF that was printed: \texttt{test\_\allowbreak{}run\_\allowbreak{}end\_\allowbreak{}writes\_\allowbreak{}real\_\allowbreak{}pdf\_\allowbreak{}report},
the two \texttt{\_selftest --run} cases and the ordinary-user-with-the-flag case.
\texttt{pdf\_has\_text(pdf)} is the check. If you add a report feature, decide which side of that line
your test is on.

\cando{add something to the Test Report and know whether your test needs a real render or a stub.}
